\documentclass[11pt,a4paper]{article}
\usepackage[margin=2.2cm]{geometry}
\usepackage[indonesian]{babel}
\usepackage{lmodern}
\usepackage[T1]{fontenc}
\usepackage{xcolor}
\usepackage{hyperref}
\usepackage{enumitem}
\usepackage{tcolorbox}
\usepackage{listings}
\usepackage{titlesec}
\usepackage{parskip}
\usepackage{fancyhdr}
\definecolor{Navy}{HTML}{141F36}
\definecolor{Blue}{HTML}{3166DC}
\definecolor{Mint}{HTML}{43BEA4}
\definecolor{Light}{HTML}{F5F7FB}
\definecolor{CodeBG}{HTML}{1C212B}
\definecolor{Gray}{HTML}{697180}
\hypersetup{colorlinks=true,linkcolor=Blue,urlcolor=Blue}
\pagestyle{fancy}\fancyhf{}\lhead{\textcolor{Navy}{DAWET}}\rhead{\textcolor{Gray}{Meeting 03}}\cfoot{\thepage}
\titleformat{\section}{\Large\bfseries\color{Navy}}{}{0em}{}
\titleformat{\subsection}{\large\bfseries\color{Blue}}{}{0em}{}
\lstdefinestyle{htmlstyle}{language=HTML,basicstyle=\ttfamily\small\color{white},backgroundcolor=\color{CodeBG},keywordstyle=\color{Mint},commentstyle=\color{Gray},stringstyle=\color{yellow},breaklines=true,frame=none,showstringspaces=false,columns=fullflexible}
\newtcolorbox{conceptbox}{colback=Light,colframe=Blue,boxrule=0.7pt,arc=3mm,left=4mm,right=4mm,top=2mm,bottom=2mm}
\begin{document}
\begin{titlepage}\centering\vspace*{2.7cm}{\Huge\bfseries\color{Navy} DAWET\par}\vspace{.5cm}{\Huge\bfseries\color{Blue} HTML Lanjutan\par}\vspace{.7cm}{\Large Meeting 03 — Modul Belajar Pemula\par}\vspace{1.4cm}\begin{tcolorbox}[colback=Light,colframe=Mint,arc=3mm]\centering Dari satu dokumen HTML menjadi website kecil yang saling terhubung.\end{tcolorbox}\vfill{\large Study Club Web Development\par}\end{titlepage}
\tableofcontents\newpage
\section{Pengantar: Dari Satu Halaman Menjadi Website}
Pada pertemuan sebelumnya, kita sudah membuat sebuah dokumen HTML yang cukup lengkap. Kita sudah menggunakan struktur dasar HTML, teks, gambar, list, tabel, serta bagian \texttt{header}, \texttt{main}, dan \texttt{footer}.
\begin{conceptbox}\textbf{Pertanyaan berikutnya:} Kalau sebuah website memiliki banyak informasi, apakah semuanya harus dimasukkan ke satu file HTML?
\end{conceptbox}
Tidak harus. Sebuah website dapat terdiri dari beberapa halaman HTML. Website UMKM, misalnya, dapat memiliki halaman Home, Tentang, Produk, dan Kontak. Setiap halaman dapat dibuat sebagai file HTML yang berbeda.
Agar pengunjung dapat berpindah dari satu halaman ke halaman lain, kita membutuhkan \textbf{link}.
\subsection{Mengapa Link Dipelajari Sekarang?}
Pada Week 2 kita sudah menggunakan attribute seperti \texttt{src} pada gambar:
\begin{lstlisting}[style=htmlstyle]<img src="img1.jpg" alt="Foto produk">\end{lstlisting}
\texttt{src} memberi tahu browser lokasi file yang dibutuhkan. Konsep yang mirip muncul pada \texttt{href}. Pada element \texttt{<a>}, \texttt{href} menentukan tujuan link.
\begin{lstlisting}[style=htmlstyle]<a href="tentang.html">Tentang Kami</a>\end{lstlisting}
Jadi kita sedang memperluas konsep yang sudah dikenal: \textbf{attribute dan path}.
\subsection{Dari File ke Website}
Misalnya:
\begin{lstlisting}[style=htmlstyle]
portal-umkm/
├── index.html
├── tentang.html
├── produk.html
└── kontak.
\end{lstlisting}
File-file tersebut perlu dihubungkan. Kumpulan link dapat menjadi navigasi:
\begin{lstlisting}[style=htmlstyle]
<a href="index.html">Home</a>
<a href="tentang.html">Tentang</a>
<a href="produk.html">Produk</a>
<a href="kontak.html">Kontak</a>
\end{lstlisting}
\subsection{Mengapa Semantic HTML Ikut Dibahas?}
Ketika halaman dan konten bertambah, struktur HTML perlu tetap jelas. Kita sudah menggunakan \texttt{header}, \texttt{main}, dan \texttt{footer}; sekarang kita menambahkan \texttt{nav}, \texttt{section}, dan pengenalan singkat \texttt{article}.
Tujuannya bukan membuat halaman langsung cantik. Tampilan akan menjadi fokus saat CSS dipelajari. Tujuan sekarang adalah membuat HTML dengan struktur dan makna yang jelas.

\section{Tujuan Pembelajaran}
\begin{enumerate}[label=\arabic*.]
\item Menjelaskan fungsi hyperlink.
\item Membuat link ke website lain.
\item Membuat link antar halaman HTML.
\item Memahami relative path sederhana.
\item Membuat navigasi dengan \texttt{nav}.
\item Memahami dasar \texttt{section} dan \texttt{article}.
\item Membuat website sederhana multi-page.
\end{enumerate}

\section{Hyperlink dengan \texttt{<a>}}
\begin{lstlisting}[style=htmlstyle]
<a href="https://www.instagram.com/">Instagram</a>
\end{lstlisting}
\texttt{href} menentukan tujuan link. Teks di antara tag pembuka dan penutup menjadi bagian yang dapat diklik.
\subsection{Link Eksternal}
\begin{lstlisting}[style=htmlstyle]
<a href="https://www.instagram.com/">Instagram</a>
<a href="https://maps.google.com/">Google Maps</a>
\end{lstlisting}
Link eksternal mengarah ke website lain dan biasanya menggunakan URL lengkap seperti \texttt{https://...}.

\section{Link Antar Halaman}
Misalnya:
\begin{lstlisting}[style=htmlstyle]
portal-umkm/
├── index.html
├── tentang.html
├── produk.html
└── kontak.html
\end{lstlisting}
Dari \texttt{index.html}:
\begin{lstlisting}[style=htmlstyle]
<a href="tentang.html">Tentang Kami</a>
<a href="produk.html">Produk</a>
<a href="kontak.html">Kontak</a>
\end{lstlisting}
Browser mencari file berdasarkan path relatif terhadap file HTML saat ini.

\section{Relative Path}
Konsep path pada link mirip dengan path gambar:
\begin{lstlisting}[style=htmlstyle]
<img src="images/menu.jpg" alt="Menu">
<a href="pages/menu.html">Menu</a>
\end{lstlisting}
\begin{conceptbox}\textbf{Cara berpikir:} Dari file HTML ini, saya harus pergi ke file mana? Tuliskan path file tersebut pada \texttt{href}.
\end{conceptbox}

\section{Navigasi dengan \texttt{<nav>}}
\begin{lstlisting}[style=htmlstyle]
<nav>
    <a href="index.html">Home</a>
    <a href="tentang.html">Tentang</a>
    <a href="produk.html">Produk</a>
    <a href="kontak.html">Kontak</a>
</nav>
\end{lstlisting}
\texttt{nav} memberi makna bahwa kumpulan link tersebut merupakan navigasi. Tampilan visualnya akan kita atur dengan CSS.

\section{Semantic HTML}
Semantic HTML berarti menggunakan element berdasarkan \textbf{makna atau perannya}, bukan sekadar berdasarkan tampilannya.
\subsection{\texttt{<header>}}Bagian pengantar atau kepala halaman/bagian.
\subsection{\texttt{<nav>}}Bagian navigasi.
\subsection{\texttt{<main>}}Konten utama halaman.
\subsection{\texttt{<section>}}Kelompok konten dengan tema atau bagian tertentu.
\subsection{\texttt{<article>}}Konten yang dapat dipandang sebagai unit tersendiri, misalnya artikel atau posting.
\subsection{\texttt{<footer>}}Bagian penutup, misalnya copyright atau informasi tambahan.

\section{Section dalam Halaman}
\begin{lstlisting}[style=htmlstyle]
<main>
    <section>
        <h2>Profil UMKM</h2>
        <p>...</p>
    </section>
    <section>
        <h2>Menu Produk</h2>
        <p>...</p>
    </section>
    <section>
        <h2>Promo</h2>
        <p>...</p>
    </section>
</main>
\end{lstlisting}

\section{Membangun Website Multi-page}
Target:
\begin{lstlisting}[style=htmlstyle]
portal-umkm/
├── index.html
├── tentang.html
├── produk.html
├── kontak.html
└── images/
    └── logo.jpg
\end{lstlisting}
Contoh Home:
\begin{lstlisting}[style=htmlstyle]
<!DOCTYPE html>
<html lang="id">
<head>
    <meta charset="UTF-8">
    <meta name="viewport" content="width=device-width, initial-scale=1.0">
    <title>Warung Berkah</title>
</head>
<body>
<header>
    <h1>Warung Berkah</h1>
    <nav>
        <a href="index.html">Home</a>
        <a href="tentang.html">Tentang</a>
        <a href="produk.html">Produk</a>
        <a href="kontak.html">Kontak</a>
    </nav>
</header>
<main>
    <section>
        <h2>Selamat Datang</h2>
        <p>Makanan rumahan untuk mahasiswa.</p>
    </section>
</main>
<footer>
    <p>&copy; 2026 Warung Berkah</p>
</footer>
</body>
</html>
\end{lstlisting}

\section{Praktik Mandiri}
Kembangkan halaman Week 2 menjadi website kecil. Minimal buat Home, Tentang, dan Produk.
\subsection{Ketentuan}
\begin{itemize}
\item Minimal tiga file HTML.
\item Struktur dasar pada setiap halaman.
\item Navigasi yang menghubungkan halaman.
\item Minimal satu link eksternal.
\item Penggunaan \texttt{header}, \texttt{nav}, \texttt{main}, \texttt{section}, dan \texttt{footer} secara masuk akal.
\item Minimal satu gambar.
\end{itemize}
\subsection{Bonus}
Tambahkan halaman Kontak, Google Maps, Instagram/WhatsApp, folder gambar yang rapi, dan \texttt{article} untuk konten mandiri.

\section{Checkpoint Pemahaman}
\begin{enumerate}
\item Apa perbedaan link eksternal dan internal?
\item Apa fungsi \texttt{href}?
\item Mengapa path \texttt{href} mirip dengan \texttt{src}?
\item Mengapa kumpulan link navigasi dibungkus dengan \texttt{nav}?
\item Apa fungsi \texttt{section}?
\item Kapan \texttt{article} masuk akal digunakan?
\item Apa yang dimaksud semantic HTML?
\end{enumerate}

\section{Tugas GitHub}
Simpan:
\begin{lstlisting}[style=htmlstyle]
DAWET/
└── minggu-03/
    └── portal-umkm/
        ├── index.html
        ├── tentang.html
        ├── produk.html
        ├── kontak.html
        └── images/
\end{lstlisting}
Lakukan commit dan push. README berisi struktur halaman, link, semantic element, kendala, hal baru, dan screenshot.

\section{Persiapan Menuju CSS}
Pada tahap ini website mungkin masih terlihat sederhana. Itu normal. HTML terutama menyusun struktur dan makna konten. CSS nantinya mengatur warna, ukuran teks, jarak, border, layout, dan tampilan halaman.
\begin{conceptbox}\textbf{Jangan terlalu sibuk mempercantik HTML.} Kita akan belajar cara yang lebih tepat melalui CSS.
\end{conceptbox}

\section{Referensi}
\begin{enumerate}
\item MDN — Anchor Element: \url{https://developer.mozilla.org/en-US/docs/Web/HTML/Reference/Elements/a}
\item MDN — nav: \url{https://developer.mozilla.org/en-US/docs/Web/HTML/Reference/Elements/nav}
\item MDN — Semantics: \url{https://developer.mozilla.org/en-US/docs/Glossary/Semantics}
\end{enumerate}

\section{Penutup}
Week 2: membuat dokumen HTML. Week 3: membuat sekumpulan halaman yang saling terhubung.
\begin{center}\textbf{Elemen HTML $\rightarrow$ Struktur halaman $\rightarrow$ Link $\rightarrow$ Multi-page website}\end{center}
\begin{conceptbox}\centering\textbf{HTML memberi kita struktur.}\\CSS akan memberi kita tampilan.\end{conceptbox}
\end{document}
